\chapter{Recapitulare structuri algebrice și matrice}

\section{Structuri algebrice}

Amintim cîteva noțiuni esențiale în ce privește structurile algebrice. Începem
cu noțiuni preliminare, anume \emph{relații binare}.

\begin{definition}\label{def:rel-bin} \index{relație!binară}
  Fie $ A $ o mulțime nevidă. O submulțime a produsului cartezian $ A \times A $
  se numește \emph{relație binară pe $ A $}.
\end{definition}

Dacă $ R \seq A \times A $ este o asemenea relație, în locul notației
$ a, b \in R $, pentru două elemente aflate în relație, vom folosi $ a R b $ și
vom citi \qq{$a$ este în relația $ R $ cu b}.

De exemplu, relația de rudenie este o relație binară, definită pe mulțimea oamenilor,
relația de ordine este o relație binară definită pe mulțimea numerelor naturale, de
pildă, relația de paralelism este o relație binară definită pe mulțimea curbelor din
plan și altele.

Relațiile binare pot avea una sau mai multe din următoarele proprietăți:

Vom păstra notațiile și contextul de mai sus. Relația $R$ se numește:
\begin{itemize} \index{relație!reflexivă} \index{relație!simetrică}
  \index{relație!antisimetrică} \index{relație!tranzitivă}
  \item \emph{reflexivă}, dacă orice element este în relația $R$ cu sine, i.e.\
    $a R a, \forall a \in A$;
  \item \emph{simetrică}, dacă oricînd $a R b$ implică $b R a$, pentru niște $a, b \in A$;
  \item \emph{antisimetrică}, dacă din $a R b$ și $b R a$, pentru niște $a, b \in A$,  
    rezultă că $a$ și $b$ coincid.
  \item \emph{tranzitivă}, dacă din $a R b$ și $b R c$ putem deduce $a R c$, pentru orice
    $a, b, c \in A$ care satisfac ipoteza.
\end{itemize}
\index{relație!ordine} \index{relație!echivalență}
În funcție de proprietățile respectate, relațiile se pot clasifica. De exemplu, cel mai des
se întîlnește relația de \emph{ordine}, care este reflexivă, antisimetrică și tranzitivă,
precum și relația de \emph{echivalență}, care este reflexivă, simetrică și tranzitivă.

Exemple clasice: relația \qq{$\leq$} este o relație de ordine pe mulțimea $\RR$. Egalitatea
este o relație de echivalență, definită pe aproape orice mulțime de obiecte matematice de
același tip (mulțimi de numere, de matrice, de funcții ș.a.m.d.).

\vspace{1cm}

Dacă relațiile binare nu returnează un al treilea element, deci nu dau o dependență funcțională,
ele avînd doar o valoare de adevăr (i.e.\ putem spune doar dacă $aRb$ este adevărat, nu și
\qq{cît face} $a R b$), \emph{operațiile} sau \emph{legile de compoziție} au această proprietate.
\index{operație}
Mai precis:

\begin{definition}\label{def:lege-comp}
  Fie $A$ o mulțime nevidă. Se numește \emph{lege de compoziție} sau \emph{operație (binară)} orice
  funcție de tipul $f : A \times A \to A$.
\end{definition}

Exemplele tipice sînt operațiile algebrice, pe mulțimi de numere. Însă înainte de a le folosi,
mai avem nevoie de terminologie.

Păstrînd notațiile și contextul de mai sus, o lege de compoziție $f$ se numește:
\begin{itemize} \index{operație!internă} \index{operație!asociativă}
  \index{operație!comutativă} \index{operație!cu element neutru}
  \index{operație!cu elemente inversabile}
  \item \emph{internă}, dacă $\forall a, b \in A, f(a, b) \in A$ (automat verificată, dacă
    în definiție impunem condiția ca $f$ să fie \emph{funcție} și nu doar o asociere abstractă);
  \item \emph{asociativă}, dacă $\forall a, b, c \in A, f(a, f(b, c)) = f(f(a, b), c)$;
  \item \emph{comutativă}, dacă $\forall a, b \in A$, are loc $f(a, b) = f(b, a)$;
  \item \emph{cu element neutru}, dacă există un element distins (unic) $e \in A$, astfel încît
    $f(a, e) = f(e, a) = a, \forall a \in A$;
  \item \emph{cu elemente inversabile (simetrizabile)} dacă pentru un $a \in A$ există $a^{-1} \in A$
    astfel încît $f(a, a^{-1}) = f(a^{-1}, a) = e$, elementul neutru.
\end{itemize}

Cum spuneam, exemplele tipice sînt operațiile algebrice elementare. Cum ar fi adunarea, definită
pe mulțimea numerelor reale, care are toate proprietățile de mai sus. Înmulțirea, definită pe
mulțimea numerelor naturale, nu admite decît un singur element inversabil, anume 1.

Compunerea funcțiilor reale, pe de altă parte, nu este comutativă, dar are elemente inversabile
(chiar dacă nu toate), anume funcțiile bijective. 

Există suficiente exemple de operații care nu au una sau mai multe din proprietățile de mai sus,
după cum vă puteți convinge.

\vspace{1cm}

Deoarece, ca în cazul relațiilor binare, există o legătură strînsă între operație și mulțimea
pe care a fost definită, se definesc \emph{structurile algebrice} ca fiind exact perechi
alcătuite din mulțimi și operații definite pe ele, satisfăcînd anumite proprietăți,
conform definițiilor de mai jos.

\begin{definition}\label{def:str-alg} \index{structuri algebrice!monoid}
  \index{structuri algebrice!grup}
  Fie $A$ o mulțime nevidă și $\circ$ o operație definită pe $A$.

  Perechea $(A, \circ)$ se numește:
  \begin{itemize}
    \item \emph{monoid}, dacă operația este internă, asociativă și admite element neutru;
    \item \emph{grup}, dacă este monoid și orice element este inversabil.
  \end{itemize}

  Dacă, în plus, operația este și comutativă, se adaugă adjectivul \qq{comutativ} pentru
  structura corespunzătoare (i.e.\ monoid comutativ, grup comutativ). În onoarea
  matematicianului norvegian N. H. Abel, grupurile comutative se mai numesc \emph{abeliene}.
\end{definition}

Există situații în care ar fi de folos să putem lucra cu mai multe operații pe o mulțime.
Gîndiți-vă, de exemplu, la cazul numerelor reale, pe care le putem și aduna, și înmulți, ba
chiar putem opera în expresii care să conțină ambele legi, simultan.
Pentru aceasta, se pot defini structuri algebrice mai bogate, adăugînd condiții suplimentare
de compatibilitate pentru operațiile folosite.

\begin{definition}\label{def:inel-corp}
  \index{structuri algebrice!inel}
  \index{structuri algebrice!corp}
  Fie $A$ o mulțime nevidă și $\circ, \ast$ două operații, ambele definite pe $A$.

  Tripletul $(A, \circ, \ast)$ se numește:
  \begin{itemize} \index{operație!distributivă}
    \item \emph{inel}, dacă perechea $(A, \circ)$ este grup comutativ, perechea $(A, \ast)$
      este monoid, iar operația $\ast$ este \emph{distributivă} față de $\circ$, adică au loc:
      \[
        a \ast (b \circ c) = (a \ast b) \circ (a \ast c) \text{ și } %
        (b \circ c) \ast a = (b \ast a) \circ (c \ast a), \forall a, b, c \in A.
      \]
    \item \emph{corp}, dacă este inel și, în plus, perechea $(A, \ast)$ este chiar grup.
  \end{itemize}

  Ca mai sus, dacă și operația $\ast$ este comutativă, se adaugă atributul \qq{comutativ}
  numelui structurii (i.e.\ inel comutativ, corp comutativ).
\end{definition}

Exemplele tipice sînt la îndemînă: $(\RR, +, \cdot)$ este corp comutativ,
$(M_n(\RR), +, \cdot)$ este inel necomutativ, $(\RR[X], +, \cdot)$ este inel
comutativ și altele.

Făcînd verificările, veți recunoaște în proprietatea de distributivitate
procedura obișnuită de \emph{desfacere a parantezelor} într-o expresie algebrică.

Date două structuri algebrice, eventual cu mulțimi subiacente și operații diferite,
există o modalitate de a le relaționa, prin intermediul \emph{morfismelor}:

\begin{definition}\label{def:morfism} \index{structuri algebrice!morfism}
  Fie $(G, \circ)$ și $(H, \ast)$ două structuri algebrice (monoizi, grupuri).

  O funcție $f: G \to H$ se numește \emph{morfism} dacă respectă operațiile, i.e.:
  \[
    f(g_1 \circ g_2) = f(g_1) \ast f(g_2), \forall g_1, g_2 \in G.
  \]

  În plus, cerem ca morfismul să fie și \emph{unitar}, adică $f(e_G) = e_H$
  (elementele neutre corespund).
\end{definition}

Dacă vrem să definim morfismele pentru inele, cerem în plus ca funcțiile să
respecte și cea de-a doua operație, iar ambele elemente neutre să corespundă.

Dacă morfismul este chiar o funcție bijectivă, atunci el se numește \emph{izomorfism},
iar structurile, \emph{izomorfe}.

O ultimă noțiune utilă:

\begin{definition}\label{def:substruct} \index{structuri algebrice!substructuri}
  Dată o structură algebrică $(G, \circ)$ și o submulțime a sa, $H$, spunem
  că $H$ este \emph{substructură} pentru $G$ (în particular, submonoid,
  subgrup) dacă $(H, \circ)$ are aceeași structură ca $G$, i.e.\ de monoid
  sau grup.
\end{definition}

În cazul structurilor cu două operații, definiția se păstrează, cerînd ca
$H$ să aibă aceeași structură cu $G$, împreună cu \emph{ambele} operații.

%%%%%%%%%%%%%%%%%%%%%%%%%%%%%%%%%%%%%%%%%%%%%%%%%%%%%%%%%%%%%%%%%%%%%%

\section{Metoda Gauss-Jordan}
\index{metoda Gauss-Jordan}
Această metodă ne permite să aducem orice matrice la o formă foarte simplă,
anume forma matricei identitate, operînd numai cu transformări elementare
asupra liniilor sau coloanelor acesteia.

Într-o formă mai relaxată, numită de obicei \emph{eliminarea gaussiană}, este suficient
să aducem matricea în formă \emph{superior triunghiulară}, adică astfel încît
elementele de sub diagonala principală să fie toate nule.

Metoda Gauss-Jordan are 2 principale aplicații (pe lîngă simplificarea matricei,
care ar putea să fie utilă în numeroase contexte). Prima aplicație este în aflarea
inversei unei matrice, iar cea de-a doua, în rezolvarea sistemelor liniare.

\subsection{Inversarea matricelor cu metoda Gauss-Jordan}

Procedura se bazează pe următorul algoritm: \index{metoda Gauss-Jordan!inversare}
\begin{itemize}
  \item Dată matricea $A$, pătrată și nesingulară (inversabilă), se alcătuiește
    matricea $B = (A \mid I_n)$, unde $n$ este dimensiunea matricei. Matricea $B$
    va fi o matrice cu $n$ linii și $2n$ coloane, în care trasăm, pentru conveniență,
    linia verticală care separă cele două componente;
  \item În partea stîngă a matricei $B$ (corespunzătoare matricei $A$), operăm cu
    \emph{transformări elementare} ale liniilor și coloanelor, pînă cînd această parte
    devine $I_n$, transformările reflectîndu-se și în partea dreaptă
    a matricei $B$;
  \item Atunci cînd în partea stîngă $A$ a devenit $I_n$, partea dreaptă a devenit $A^{-1}$.
\end{itemize}

Iată un exemplu, notînd și transformările:
\begin{align*}
  \begin{pmatrix}
    2 & 1 & 3 & \mid & 1 & 0 & 0 \\
    -2 & 3 & 4 & \mid & 0 & 1 & 0 \\
    5 & 1 & 1 & \mid & 0 & 0 & 1
  \end{pmatrix} &\xrightarrow{L_1 \to (1/2)L_1}
  \begin{pmatrix}
    1 & 1/2 & 3/2 & \mid & 1/2 & 0 & 0 \\
    -2 & 3 & 4 & \mid & 0 & 1 & 0 \\
    5 & 1 & 1 & \mid & 0 & 0 & 1 
  \end{pmatrix} \\
  &\stackrel{L_2 \to 2L_1 + L_2}{\xrightarrow{L_3 \to -5L_1 + L_3}}
  \begin{pmatrix}
    1 & 1/2 & 3/2 & \mid & 1/2 & 0 & 0 \\
    0 & 4 & 7 & \mid & 1 & 1 & 0 \\
    0 & -3/2 & -13/2 & \mid & -5/2 & 0 & 1
  \end{pmatrix} \\
  &\xrightarrow{L_2 \to (1/4)L_2}
  \begin{pmatrix}
    1 & 1/2 & 3/2 & \mid & 1/2 & 0 & 0 \\
    0 & 1 & 7/4 & \mid & 1/4 & 1/4 & 0 \\
    0 & -3/2 & -13/2 & \mid & -5/2 & 0 & 1 
  \end{pmatrix} \\
  &\stackrel{L_1 \to (-1/2)L_2 + L_1}{\xrightarrow{L_3 \to (3/2)L_2 + L_3}}
  \begin{pmatrix}
    1 & 0 & 5/8 & \mid & 3/8 & -1/8 & 0 \\
    0 & 1 & 7/4 & \mid & 1/4 & 1/4 & 0 \\
    0 & 0 & -31/8 & \mid & -17/8 & 3/8 & 1
  \end{pmatrix} \\
  &\xrightarrow{L_3 \to (-31/8)L_3}
  \begin{pmatrix}
    1 & 0 & 5/8 & \mid & 3/8 & -1/8 & 0 \\
    0 & 1 & 7/4 & \mid & 1/4 & 1/4 & 0 \\
    0 & 0 & 1 & \mid & 17/31 & -3/31 & -8/31
  \end{pmatrix} \\
  &\stackrel{L_1 \to (-5/8)L_3 + L_1}{\xrightarrow{L_2 \to (-7/4)L_3 + L_2}}
  \begin{pmatrix}
    1 & 0 & 0 & \mid & 1/31 & -2/31 & 5/31 \\
    0 & 1 & 0 & \mid & -22/31 & 13/31 & 14/31 \\
    0 & 0 & 1 & \mid & 17/31 & -3/31 & -8/31
  \end{pmatrix}
\end{align*}

Rezultă că $A^{-1} = \begin{pmatrix} 1/31 & -2/31 & 5/31 \\
  -22/31 & 13/31 & 14/31 \\
  17/31 & -3/31 & -8/31
\end{pmatrix}$.

%%%%%%%%%%%%%%%%%%%%%%%%%%%%%%%%%%%%%%%%%%%%%%%%%%%%%%%%%%%%%%%%%%%%%%

\subsection{Aplicație: Sisteme liniare}
\index{metoda Gauss-Jordan!sisteme liniare}
Vom proceda similar, însă vom opera pe matricea $M = (A | B)$, unde
$A$ este matricea asociată sistemului liniar, iar $B$ este matricea
termenilor liberi (rezultatelor). Cînd, după procedeul Gauss-Jordan,
am obținut $M = (I_n | C)$, atunci $C$ va fi matricea-coloană a soluțiilor.

Iată un exemplu: Considerăm sistemul de mai jos.
\[
  \begin{cases}
    mx_1 + x_2 + x_3  &= 1 \\
    x_1 + mx_2 + x_3 &= m \\
    x_1 + x_2 + mx_3 &= m^2,
  \end{cases}
\]
cu $m \in \RR$. Îi vom discuta compatibilitatea în funcție de $m$ și, în caz
favorabil, îl vom rezolva. Totul va rezulta din metoda Gauss-Jordan.

Considerăm matricea $M$ de mai sus:
\begin{align*}
  (A | B) &= \begin{pmatrix}
    m & 1 & 1 & \mid & 1 \\
    1 & m & 1 & \mid & m \\
    1 & 1 & m & \mid & m^2
  \end{pmatrix} \\
  &\xrightarrow{L_1 \leftrightarrow L_2}
  \begin{pmatrix}
    1 & m & 1 & \mid & m \\
    m & 1 & 1 & \mid & 1 \\
    1 & 1 & m & \mid & m^2
  \end{pmatrix} \\
  &\stackrel{L_2 \to L_2 - mL_1}{\xrightarrow{L_3 \to L_3 - L_1}}
  \begin{pmatrix}
    1 & m & 1 & \mid & m \\
    0 & 1 - m^2 & 1 - m & \mid & 1 - m^2 \\
    0 & 1 - m & m - 1 & \mid & m^2 - m.
  \end{pmatrix}
\end{align*}

În acest punct, avem o discuție:

(a) Dacă $m = 1$, atunci sistemul se reduce la prima ecuație, deci este compatibil
dublu nedeterminat. Rezultă soluția
\[
  \{ (1 - \alpha - \beta, \alpha, \beta) \mid \alpha, \beta \in \RR \}.
\]

(b) Dacă $m \neq 1$, atunci putem continua transformările și ajungem, în fine, la:
\[
  M = \begin{pmatrix}
    1 & 0 & 1 + m & \mid & m + m^2 \\
    0 & 1 & -1 & \mid & -m \\
    0 & 0 & 2 & \mid & (m + 1)^2
  \end{pmatrix}
\]

Aici discutăm din nou:

(b1) Dacă $m = -2$, sistemul este incompatibil, deoarece ultima ecuație devine $0 = 1$.

(b2) Dacă $m \neq -2$, putem continua transformările și ajungem, în cele din urmă, la:
\[
  M = \begin{pmatrix}
    1 & 0 & 0 & \mid & - \frac{m+1}{m_2}  \\
    0 & 1 & 0 & \mid & \frac{1}{m+2} \\
    0 & 0 & 1 & \mid & \frac{(m + 1)^2}{m + 2}.
  \end{pmatrix}
\]

În acest ultim caz, sistemul este compatibil determinat, soluția fiind dată de ultima coloană:
\[
  \begin{cases}
    x_1  &= -\frac{m + 1}{2} \\
    x_2  &= \frac{1}{m + 2} \\
    x_3  &= \frac{(m + 1)^2}{m + 2}.
  \end{cases}
\]

Concluzia generală este:
\begin{enumerate}[(a)]
  \item Dacă $m \in \RR - \{-2, 1\}$, sistemul are soluție unică;
  \item Dacă $m = -2$, sistemul este incompatibil;
  \item Dacă $m = -1$, sistemul este compatibil nedeterminat.
\end{enumerate}

%%%%%%%%%%%%%%%%%%%%%%%%%%%%%%%%%%%%%%%%%%%%%%%%%%%%%%%%%%%%%%%%%%%%%%

\section{Exerciții}

1. Rezolvați următoarele sisteme, atît cu metoda matriceală clasică, cît și cu
metoda lui Gauss:
\begin{enumerate}[(a)]
\item $ \begin{cases}
    x + 3y - z &= 1 \\
    3x + y + 2z &= 4 \\
    5x - 2y + z &= 2
  \end{cases} $
\item $ \begin{cases}
    3x - y + 2z + 3t &= 4 \\
    x + 3y - z + t &= -1
  \end{cases} $
\item $ \begin{cases}
    2x - y + 3z &= 1 \\
    3x + 2y - z &= 4 \\
    x + 3y - 4z &= 3
  \end{cases} $.
\end{enumerate}

\vspace{1cm}

2. Calculați inversele matricelor, atît cu metoda folosind matricea adjunctă,
cît și folosind metoda lui Gauss:
\begin{enumerate}[(a)]
\item $ A =
  \begin{pmatrix}
1 & 0 & -1 \\
3 & 2 & 1 \\
0 & 1 & 1
\end{pmatrix} $;
\item $ B =
  \begin{pmatrix}
    -1 & 2 & 4 & 1 \\
    3 & 0 & 1 & 4 \\
    0 & 1 & 1 & 0 \\
    2 & 1 & 0 & 1
  \end{pmatrix} $;
\item $ C =
  \begin{pmatrix}
    1 & 1 & 2 \\
    0 & 1 & 0 \\
    0 & 0 & -1
  \end{pmatrix} $.
\end{enumerate}

%%% Local Variables:
%%% mode: latex
%%% TeX-master: "../ag-18-19"
%%% End:
